% PDF-to-TeX transcription by the assistant; not the original downloaded file.
\begin{definition}[5.17]
Let $E$ be a finite Galois extension of $F$. A basis for $E$ as an $F$-vector space is called a normal basis if it consists of the conjugates of a single element of $E$. In other words, a normal basis is one of the form
\[
\{\sigma\alpha\mid\sigma\in\operatorname{Gal}(E/F)\}
\]
for some $\alpha\in E$.
\end{definition}
\begin{theorem}[5.18, Normal basis theorem]
Every Galois extension has a normal basis.
\end{theorem}
The group algebra $FG$ of a group $G$ is the $F$-vector space with basis the elements of $G$ endowed with the multiplication extending that of $G$. Every $F$-linear action of $G$ on an $F$-vector space $V$ extends uniquely to an action of $FG$ on $V$.

Let $E/F$ be a Galois extension with Galois group $G$. Then $E$ is an $FG$-module, and Theorem 5.18 says that there exists an element $\alpha\in E$ such that the map
\[
\sum_{\sigma}a_\sigma\sigma\longmapsto\sum_{\sigma}a_\sigma\sigma\alpha:FG\longrightarrow E
\]
is an isomorphism of $FG$-modules, i.e., that $E$ is a free $FG$-module of rank $1$.

\begin{proof}[Uniform proof]
Recall that a module is indecomposable if it is nonzero and cannot be written as a direct sum of two nonzero submodules. The Krull--Schmidt theorem says that every nonzero module $M$ of finite length over a ring can be written as a direct sum of indecomposable modules and that the indecomposable modules occurring in a decomposition are unique up to order and isomorphism. Thus $M=\bigoplus_i m_iM_i$ where $M_i$ is indecomposable and $m_iM_i$ denotes the direct sum of $m_i$ copies of $M_i$; the set of isomorphism classes of the $M_i$ is uniquely determined and, when we choose the $M_i$ to be pairwise nonisomorphic, each $m_i$ is uniquely determined. From this it follows that two modules $M$ and $M'$ of finite length over a ring are isomorphic if $mM\simeq mM'$ for some $m\geq1$.

Consider the $F$-vector space $E\otimes_F E$. We let $E$ act on the first factor, and $G$ act on the second factor (so $a(x\otimes y)=ax\otimes y$, $a\in E$, and $\sigma(x\otimes y)=x\otimes\sigma y$, $\sigma\in G$). We'll prove Theorem 5.18 by showing that
\[
\underbrace{FG\oplus\cdots\oplus FG}_{n}\simeq E\otimes_F E
\simeq\underbrace{E\oplus\cdots\oplus E}_{n}
\]
as $FG$-modules ($n=[E:F]$).

For $\sigma\in G$, let $\lambda_\sigma:E\otimes_F E\to E$ denote the map $x\otimes y\mapsto x\cdot\sigma y$. Then $\lambda_\sigma$ is obviously $E$-linear, and $\lambda_\sigma(\tau z)=\lambda_{\sigma\tau}(z)$ for all $\tau\in G$ and $z\in E\otimes_F E$. I claim that $\{\lambda_\sigma\mid\sigma\in G\}$ is an $E$-basis for $\Hom_{E\text{-linear}}(E\otimes_F E,E)$. As this space has dimension $n$, it suffices to show that the set is linearly independent. But if $\sum_\sigma c_\sigma\lambda_\sigma=0$, $c_\sigma\in E$, then
\[
0=\sum_\sigma c_\sigma\bigl(\lambda_\sigma(1\otimes y)\bigr)
=\sum_\sigma c_\sigma\cdot\sigma y
\]
for all $y\in E$, which implies that all $c_\sigma=0$ by Dedekind's theorem 5.14.

Consider the map
\[
\Phi:E\otimes_F E\longrightarrow EG,\qquad
z\longmapsto\sum_\sigma\lambda_\sigma(z)\cdot\sigma^{-1}.
\]
Then $\Phi$ is $E$-linear. If $\Phi(z)=0$, then $\lambda_\sigma(z)=0$ for all $\sigma\in G$, and so $z=0$ in $E\otimes_F E$ (because the $\lambda_\sigma$ span the dual space). Therefore $\Phi$ is injective, and as $E\otimes_F E$ and $EG$ both have dimension $n$ over $E$, it is an isomorphism. For $\tau\in G$,
\[
\Phi(\tau z)=\sum_\sigma\lambda_\sigma(\tau z)\cdot\sigma^{-1}
=\sum_\sigma\lambda_{\sigma\tau}(z)\cdot\tau(\sigma\tau)^{-1}
=\tau\Phi(z);
\]
and so $\Phi$ is an isomorphism of $EG$-modules. Thus
\[
E\otimes_K E\simeq EG\simeq FG\oplus\cdots\oplus FG
\]
as an $FG$-module.

On the other hand, for any basis $\{e_1,\ldots,e_n\}$ for $E$ as an $F$-vector space,
\[
E\otimes_F E=(e_1\otimes E)\oplus\cdots\oplus(e_n\otimes E)
\simeq E\oplus\cdots\oplus E
\]
as $FG$-modules. This completes the proof.
\end{proof}

